\textbf{斜率优化}

常见标准式可写成
\[
dp[i]=A[i]+\min_j \{Y[j]-K[i]X[j]\}
\]
或等价的最大化形式。把每个决策 $j$ 看成点 $(X[j],Y[j])$，查询 $i$ 看成给定斜率 $K[i]$，在凸壳上找最优切点。

\textbf{三种常见实现}

\begin{itemize}
	\item \textbf{单调队列}：插入的 $X[j]$ 单调，且常见地查询斜率 $K[i]$ 也单调。最简单，常数最小。
	\item \textbf{李超线段树}：在线维护直线集合，适合 $X$ 和 $K$ 都不单调但查询点 $x$ 落在整数定义域的情形。
	\item \textbf{CDQ分治}：离线处理下标有先后约束的转移，常见标准式为
	\[
	dp[i]=c[i]+\min_{0\le j<i}(dp[j]+a[i]b[j]).
	\]
\end{itemize}

\textbf{实现步骤}

\begin{itemize}
	\item 先把转移整理成 $Y[j]-K[i]X[j]$。
	\item 明确是求最小值还是最大值，从而决定维护下凸壳还是上凸壳。
	\item 判断 $X$ 是否单调、$K$ 是否单调、是否必须离线，再选实现。
	\item 处理好相同 $X$ 的点：求最小值时只保留 $Y$ 更小者。
\end{itemize}

\textbf{易错点}

\begin{itemize}
	\item 斜率比较尽量用交叉乘，少用浮点。
	\item 单调队列和 CDQ 的判劣式、查询方向都和“求最小/最大”有关，符号不要抄反。
	\item 李超线段树要确认定义域覆盖所有查询点。
	\item CDQ 只能处理“决策下标在前、转移下标在后”的离线依赖。
\end{itemize}
